\documentclass[11pt,a4paper]{article}

\usepackage[utf8]{inputenc}
\usepackage[T1]{fontenc}
\usepackage[spanish]{babel}
\usepackage{lmodern}
\usepackage[a4paper,margin=2.5cm]{geometry}
\usepackage{graphicx}
\usepackage{booktabs}
\usepackage{amsmath}
\usepackage{enumitem}
\usepackage{float}
\usepackage{tikz}
\usepackage{pgfplots}
\pgfplotsset{compat=1.16}
\usetikzlibrary{positioning}
\usepackage[hidelinks]{hyperref}

\setlength{\parindent}{0pt}
\setlength{\parskip}{0.7em}
\emergencystretch=3em

% Nombre de un campo o archivo (admite guiones bajos).
\newcommand{\campo}[1]{\texttt{\detokenize{#1}}}

\begin{document}

% ---------------------------------------------------------------- Portada
\begin{titlepage}
  \centering
  \includegraphics[height=2.6cm]{img/logo-fai.png}
  \hfill
  \includegraphics[height=2.6cm]{img/logo-unco.png}

  \vspace{2.2cm}
  {\large Departamento Teoría de la Computación\par}
  {\large\textbf{Facultad de Informática}\par}
  {\large Universidad Nacional del Comahue\par}

  \vspace{1.8cm}
  {\huge Laboratorio de inteligencia artificial\par}
  \vspace{0.4cm}
  {\Large Trabajo Práctico Obligatorio N\textsuperscript{o} 1\par}
  \vspace{0.5cm}
  {\LARGE\textbf{Parte 3: entrenamiento de los modelos\\ y explicación de los resultados}\par}
  \vspace{0.3cm}
  {\normalsize Predicción de la duración de desvíos de ruta en flotas de transporte de personal\par}

  \vspace{1.2cm}
  \begin{tabular}{@{}lll@{}}
    \toprule
    \textbf{Apellido y Nombre} & \textbf{Legajo} & \textbf{Email} \\
    \midrule
    Reibold, Sebastian Alejandro & FAI-3854 & sebastian.reibold@est.fi.uncoma.edu.ar \\
    Fernandez Gramajo, Bautista  & FAI-3874 & bautista.fernandez@est.fi.uncoma.edu.ar \\
    Kissner, Davor               & FAI-3185 & davor.kissner@est.fi.uncoma.edu.ar \\
    \bottomrule
  \end{tabular}

  \vfill
  {\large Cursado 2026 \quad---\quad 7 de noviembre de 2026\par}
\end{titlepage}

\tableofcontents
\newpage

% ---------------------------------------------------------------- 1
\section{Introducción}

El trabajo plantea una \textbf{regresión}: en el instante $T_0$ en que aparece el primer punto fuera de
ruta de un vehículo, y usando solo la información disponible hasta ese momento, predecir cuántos
segundos va a durar el desvío. Esta parte presenta los modelos entrenados sobre el corpus, cómo se los
evalúa y qué resultados dan.

El trabajo se hizo en dos etapas. En la primera se armaron una referencia (la mediana), un árbol de
decisión y un Random Forest, y se midieron con el error absoluto medio. Después de revisar esos
resultados con los docentes y el equipo de Tecfield se decidió mejorar tres cosas: agregar una métrica
más (el RMSE), analizar qué variables pesan más en los modelos y en qué orden separa el árbol los
desvíos, y probar un modelo nuevo, XGBoost. La sección~2 resume el corpus, la sección~3
describe la primera etapa, la sección~4 las mejoras y la sección~5 compara todos los modelos. La
sección~6 reúne las conclusiones.

Todo el trabajo está en el cuaderno de marimo \campo{notebooks/01_exploracion_y_modelos.py}, con
\emph{polars} para las tablas y \emph{scikit-learn} y \emph{XGBoost} para los modelos. Los
números de este documento salen de su ejecución completa.

% ---------------------------------------------------------------- 2
\section{El corpus en breve}

La unidad de análisis es el \emph{evento} de desvío (todas las lecturas de una misma salida de ruta) y la
variable objetivo $y$ es su duración en segundos, medida con la hora del GPS. El corpus tiene
\textbf{13.632 eventos} de \textbf{556 vehículos}, con una mediana de 20,4 minutos y un percentil 90 de
115,8 minutos: la distribución tiene cola larga. Los eventos que llegan al tope de cuatro horas del
detector no se usan para entrenar.

Cada evento tiene \textbf{23 atributos}, todos conocidos en $T_0$, en tres grupos (cuadro~\ref{tab:atributos}).
Cinco son categóricos: ruta más cercana, franja horaria, tipo de vehículo, tipo de hoja y obra. No se usa
ningún dato posterior a $T_0$, para que los modelos no vean el futuro.

\begin{table}[H]
  \centering
  \begin{tabular}{@{}llp{9.3cm}@{}}
    \toprule
    \textbf{Grupo} & \textbf{Cant.} & \textbf{Atributos} \\
    \midrule
    Primer ping & 9 &
      distancia a la ruta, velocidad y si viene vacía, ruta más cercana, latitud, longitud, hora, día de
      la semana y franja \\
    Hoja de ruta & 8 &
      tipo de vehículo, tipo de hoja, obra, kilómetros y minutos de la ida, plazas, pasajeros y si la hoja
      estaba vigente en $T_0$ \\
    Rastreo previo & 6 &
      cantidad de GPS, velocidad media y última, minutos desde el último GPS, desplazamiento y rumbo,
      en los 10 minutos previos a $T_0$ \\
    \bottomrule
  \end{tabular}
  \caption{Atributos conocidos en $T_0$.}
  \label{tab:atributos}
\end{table}

La cobertura es desigual: el 42\,\% de los eventos no tiene ningún GPS en los 10 minutos previos (les
faltan los 6 atributos del rastreo previo) y 6.804 eventos no tienen tipo de hoja.

% ---------------------------------------------------------------- 3
\section{Primera etapa: modelos base}

\subsection{Por qué se predice el logaritmo de la duración}

Las duraciones son muy desparejas: la mayoría de los desvíos dura minutos y unos pocos duran horas. Si el
modelo aprende la duración tal cual, esos pocos casos pesan demasiado. Por eso los modelos aprenden
$\log(1+y)$, que achica las diferencias grandes, y al predecir se vuelve a segundos con $e^{\hat z}-1$. Se
suma 1 para poder tomar el logaritmo de una duración igual a cero.

\subsection{Referencia: la mediana}

Predice \textbf{siempre la misma duración}, la mediana de los desvíos de entrenamiento (unos 20 minutos),
sin mirar ningún atributo. Es lo que haría un operador sin modelo, y es el piso que hay que superar. Se
eligió la mediana y no el promedio porque, entre las predicciones constantes, es la que comete el menor
error absoluto medio: es la referencia constante más difícil de superar.

\subsection{Árbol de decisión}

Un árbol de decisión hace una serie de preguntas sobre el desvío, una detrás de otra. Cada respuesta lleva
el desvío por una rama distinta hasta una \emph{hoja}, y la predicción es la duración típica de los desvíos
de entrenamiento que llegaron a esa hoja. Se lo limita para que no memorice los datos (cuadro~\ref{tab:arbol}).
Su ventaja es que se puede dibujar y explicar; su límite es que es \emph{inestable}: un cambio chico en los
datos puede cambiarle toda la estructura.

\begin{table}[H]
  \centering
  \begin{tabular}{@{}lrl@{}}
    \toprule
    \textbf{Parámetro} & \textbf{Valor} & \textbf{Para qué} \\
    \midrule
    \campo{max_depth}        & 8  & como máximo 8 preguntas seguidas \\
    \campo{min_samples_leaf} & 30 & cada hoja reúne al menos 30 desvíos \\
    \campo{random_state}     & 0  & resultados repetibles \\
    \bottomrule
  \end{tabular}
  \caption{Árbol de decisión (\campo{DecisionTreeRegressor}). Entrenado con todo el corpus termina con 153 hojas.}
  \label{tab:arbol}
\end{table}

\subsection{Random Forest}

Entrena \textbf{200 árboles} y los hace distintos a propósito: cada uno aprende con una muestra de los
desvíos sorteada con reposición, de modo que ve los datos de una forma un poco diferente y se equivoca de
otra manera. La predicción final es el \textbf{promedio de los 200}, y al promediar los errores que van para
un lado se compensan con los que van para el otro. Esto corrige la inestabilidad del árbol solo. La
configuración (cuadro~\ref{tab:rf}) es una primera elección, sin ajustar.

\begin{table}[H]
  \centering
  \begin{tabular}{@{}lrl@{}}
    \toprule
    \textbf{Parámetro} & \textbf{Valor} & \textbf{Para qué} \\
    \midrule
    \campo{n_estimators}     & 200 & cantidad de árboles \\
    \campo{min_samples_leaf} & 10  & cada hoja de cada árbol reúne al menos 10 desvíos \\
    \campo{random_state}     & 0   & resultados repetibles \\
    \bottomrule
  \end{tabular}
  \caption{Random Forest (\campo{RandomForestRegressor}).}
  \label{tab:rf}
\end{table}

\subsection{Detalles técnicos}

Los atributos de texto se pasan a los modelos como códigos $0,1,2,\dots$; el árbol los trata como si
tuvieran un orden que no tienen, lo que es una limitación. Los datos faltantes (por ejemplo, el rastreo
previo de un vehículo que no mandó GPS) no se completan: los propios árboles deciden por qué rama mandar el
desvío, lo que requiere scikit-learn 1.3 o superior en el árbol y 1.4 o superior en el bosque.

\subsection{Cómo se evalúa}

Para saber si un modelo funciona hay que probarlo con desvíos que no vio al entrenar. Se usa una
\textbf{validación cruzada por vehículo} de cinco particiones: los 556 vehículos se reparten en cinco grupos,
todos los desvíos de un vehículo quedan en el mismo grupo, y se entrena con cuatro grupos para predecir el
quinto, cinco veces (cuadro~\ref{tab:cv}). Así cada uno de los 13.632 desvíos tiene una predicción hecha por
un modelo que nunca vio a su vehículo, y el error se calcula una sola vez sobre todas juntas.

\begin{table}[H]
  \centering
  \begin{tabular}{@{}lccccc@{}}
    \toprule
    & \textbf{Grupo 1} & \textbf{Grupo 2} & \textbf{Grupo 3} & \textbf{Grupo 4} & \textbf{Grupo 5} \\
    \midrule
    Vuelta 1 & prueba  & entrena & entrena & entrena & entrena \\
    Vuelta 2 & entrena & prueba  & entrena & entrena & entrena \\
    Vuelta 3 & entrena & entrena & prueba  & entrena & entrena \\
    Vuelta 4 & entrena & entrena & entrena & prueba  & entrena \\
    Vuelta 5 & entrena & entrena & entrena & entrena & prueba  \\
    \bottomrule
  \end{tabular}
  \caption{Validación cruzada por vehículo.}
  \label{tab:cv}
\end{table}

Se separa por vehículo y no al azar porque, si no, un mismo vehículo podría estar en entrenamiento y en
prueba, y el modelo aprendería sus costumbres y parecería mejor de lo que es. Lo que importa en la práctica
es cómo predice para un vehículo que nunca vio. Konovalenko et al.~\cite{konovalenko} separan los datos por
fecha, con la misma intención.

\subsection{Métricas}

El \textbf{error} de cada desvío es la diferencia, sin signo, entre la duración predicha y la real. Con esos
errores se calculan estas medidas, que conviene ver con un ejemplo inventado de cuatro desvíos
(cuadro~\ref{tab:ejemplo}).

\begin{table}[H]
  \centering
  \begin{tabular}{@{}lrrr@{}}
    \toprule
    \textbf{Desvío} & \textbf{Duró (min)} & \textbf{Predicción (min)} & \textbf{Error (min)} \\
    \midrule
    A & 10  & 15 & 5 \\
    B & 20  & 18 & 2 \\
    C & 30  & 40 & 10 \\
    D & 120 & 50 & 70 \\
    \bottomrule
  \end{tabular}
  \caption{Ejemplo inventado para explicar las métricas.}
  \label{tab:ejemplo}
\end{table}

\begin{itemize}[nosep]
  \item \textbf{Error absoluto medio (MAE).} Es el promedio de los errores: en promedio, por cuántos
  minutos se equivoca el modelo. En el ejemplo, $(5+2+10+70)/4 = 21{,}75$ minutos. Es la \textbf{métrica
  principal}.
  \item \textbf{Error mediano.} Es la mediana de los errores: la mitad de las predicciones se equivoca
  menos que este valor. En el ejemplo, $(5+10)/2 = 7{,}5$ minutos.
  \item \textbf{Mejora.} Qué parte del MAE de la referencia le saca el modelo:
  $\text{mejora} = 1 - \dfrac{\text{MAE del modelo}}{\text{MAE de la mediana}}$. Para la propia referencia
  es cero.
\end{itemize}

\subsection{Resultados de la primera etapa}

\begin{table}[H]
  \centering
  \begin{tabular}{@{}lrrr@{}}
    \toprule
    \textbf{Modelo} & \textbf{MAE (min)} & \textbf{Error mediano (min)} & \textbf{Mejora} \\
    \midrule
    Mediana (referencia) & 32,1 & 15,0 & --- \\
    Árbol de decisión    & 28,8 & 14,3 & 10\,\% \\
    Random Forest        & 27,0 & 13,0 & 16\,\% \\
    \bottomrule
  \end{tabular}
  \caption{Primera etapa: error con los 23 atributos, en validación cruzada por vehículo.}
  \label{tab:etapa1}
\end{table}

Si se predice siempre la duración típica, el error promedio es de 32,1 minutos y la mitad de las veces se
erra por menos de 15. El árbol baja el MAE a 28,8 minutos (10\,\% de mejora) y el Random Forest a
27,0 (16\,\%), con la mitad de sus predicciones a menos de 13 minutos del valor real. Las dos métricas
ordenan los modelos igual. En los tres casos el MAE es bastante mayor que el error mediano: unos pocos
desvíos muy largos, con errores grandes, suben el promedio.

\subsection{Aporte de cada grupo de atributos}

Para saber de dónde sale la mejora se entrenó el Random Forest tres veces, sumando un grupo de atributos
por vez y con la misma validación (cuadro~\ref{tab:grupos}).

\begin{table}[H]
  \centering
  \begin{tabular}{@{}lrrrr@{}}
    \toprule
    \textbf{Atributos del Random Forest} & \textbf{Cant.} & \textbf{MAE (min)} & \textbf{RMSE (min)} & \textbf{Mejora} \\
    \midrule
    Solo primer ping & 9 & 29,4 & 50,2 & 8\,\% \\
    + Hoja de ruta & 17 & 27,0 & 45,7 & 16\,\% \\
    + Rastreo previo & 23 & 27,0 & 45,8 & 16\,\% \\
    \bottomrule
  \end{tabular}
  \caption{Random Forest sumando los grupos de atributos de a uno.}
  \label{tab:grupos}
\end{table}

Solo con el primer ping el bosque ya mejora un 8\,\% a la referencia. Al sumar la hoja de ruta la mejora se
duplica. Al sumar el rastreo previo el error no cambia: ese grupo no aporta información nueva, algo
esperable si se recuerda que le falta al 42\,\% de los eventos.

% ---------------------------------------------------------------- 4
\section{Segunda etapa: mejoras}

Al presentar estos resultados, el equipo docente y Tecfield sugirieron tres mejoras: \textbf{agregar una
métrica más}, el RMSE; \textbf{analizar qué variables importan}, tanto en los modelos de bosque como en el
árbol, donde interesa saber cuáles son las primeras que separan los desvíos; y \textbf{probar un modelo
más}, de la familia de \emph{gradient boosting} (XGBoost). Se hicieron las tres.

\subsection{Una métrica más: el RMSE}

El \textbf{RMSE} (raíz del error cuadrático medio) eleva los errores al cuadrado antes de promediarlos y
saca la raíz al final, de modo que queda en minutos como el MAE:
\[
  \text{RMSE} = \sqrt{\frac{1}{n}\sum_{i=1}^{n} (\hat y_i - y_i)^2}.
\]
Al elevar al cuadrado, los errores grandes pesan mucho más. En el ejemplo del cuadro~\ref{tab:ejemplo} el
RMSE es $\sqrt{(25+4+100+4900)/4} = 35{,}5$ minutos, contra 21,75 de MAE: el desvío D, con 70 minutos de
error, domina el resultado. Por eso las dos métricas se complementan: el MAE dice cuánto se equivoca el modelo
\emph{en general} y el RMSE cuánto lo castigan los errores grandes. Para este problema, donde hay pocos
desvíos muy largos, la diferencia entre ambos es informativa. El RMSE es 1,7 veces el MAE en el mejor
modelo (45,8 contra 27,0 minutos), lo que confirma que el error se concentra en los desvíos largos.

\subsection{Modelos de gradient boosting}

A diferencia del bosque, que promedia árboles independientes, el \emph{gradient boosting} arma los árboles
\textbf{de a uno}, y cada árbol nuevo intenta corregir los errores que dejaron los anteriores. Se usan
árboles chicos y muchos pasos pequeños. XGBoost es una implementación de esta idea que acepta datos faltantes
sin completarlos. Se probó con
una configuración inicial, sin ajustar (cuadro~\ref{tab:boost}), y con las mismas cinco particiones por
vehículo.

\begin{table}[H]
  \centering
  \begin{tabular}{@{}lr@{}}
    \toprule
    \textbf{Parámetro} & \textbf{Valor} \\
    \midrule
    cantidad de árboles (\campo{n_estimators}) & 300 \\
    paso de aprendizaje (\campo{learning_rate}) & 0,05 \\
    tamaño del árbol & \campo{max_depth}=4 \\
    mínimo por hoja & \campo{min_child_weight}=10 \\
    fracción de desvíos por árbol (\campo{subsample}) & 0,8 \\
    \campo{random_state} & 0 \\
    \bottomrule
  \end{tabular}
  \caption{Configuración de XGBoost.}
  \label{tab:boost}
\end{table}

\subsection{Importancia de las variables}

Cada modelo de árboles permite medir la \emph{importancia} de cada atributo: cuánto reduce el error, sumando
todas las preguntas que se hacen sobre él. Las importancias de cada modelo suman 1. Se calcularon entrenando
cada modelo una vez con el corpus completo (la figura~\ref{fig:imp} muestra los atributos más importantes).
Son importancias \emph{del modelo}, no de la realidad: dicen qué usó el modelo para predecir, no qué causa
la duración.

\begin{figure}[H]
  \centering
  \begin{tikzpicture}
    \begin{axis}[
      xbar, width=0.82\textwidth, height=10.5cm,
      bar width=3.2pt, enlarge y limits=0.06,
      xlabel={importancia}, xmin=0,
      symbolic y coords={Tipo de hoja,Distancia a la ruta,Latitud,Velocidad,Longitud,Hora,Obra,Kilómetros de ida,Minutos de ida,Pasajeros,Hoja vigente},
      ytick=data, y dir=reverse,
      legend style={at={(0.5,-0.12)}, anchor=north, legend columns=3, font=\footnotesize},
      scaled x ticks=false, xtick distance=0.1,
      tick label style={font=\footnotesize},
    ]
      \addplot[fill=blue!70,draw=blue!70] coordinates {(0.5082,Tipo de hoja) (0.1399,Distancia a la ruta) (0.0678,Latitud) (0.0631,Velocidad) (0.0384,Longitud) (0.0440,Hora) (0.0238,Obra) (0.0166,Kilómetros de ida) (0.0146,Minutos de ida) (0.0179,Pasajeros) (0.0167,Hoja vigente)};
      \addplot[fill=orange!80,draw=orange!80] coordinates {(0.2920,Tipo de hoja) (0.1255,Distancia a la ruta) (0.0902,Latitud) (0.0776,Velocidad) (0.0810,Longitud) (0.0449,Hora) (0.0357,Obra) (0.0452,Kilómetros de ida) (0.0452,Minutos de ida) (0.0256,Pasajeros) (0.0108,Hoja vigente)};
      \addplot[fill=green!50!black,draw=green!50!black] coordinates {(0.4918,Tipo de hoja) (0.0473,Distancia a la ruta) (0.0303,Latitud) (0.0314,Velocidad) (0.0250,Longitud) (0.0328,Hora) (0.0288,Obra) (0.0198,Kilómetros de ida) (0.0210,Minutos de ida) (0.0279,Pasajeros) (0.0414,Hoja vigente)};
      \legend{Árbol, Random Forest, XGBoost}
    \end{axis}
  \end{tikzpicture}
  \caption{Atributos más importantes de cada modelo (la unión de los ocho primeros de cada uno).}
  \label{fig:imp}
\end{figure}

Los tres modelos coinciden en lo principal:

\begin{itemize}[nosep]
  \item El \textbf{tipo de hoja} es el atributo más importante en todos, con entre el 29\,\% (Random Forest) y
  el 51\,\% (árbol) del total.
  \item Le siguen la \textbf{distancia a la ruta} y la \textbf{posición} (latitud y longitud) y la
  \textbf{velocidad} del primer ping, además de la hora del desvío.
  \item Los atributos del \textbf{rastreo previo} pesan poco.
\end{itemize}

Sumando por grupo (cuadro~\ref{tab:imp-grupo}), la hoja de ruta y el primer ping explican casi todo, y el
rastreo previo no supera el 10\,\% en ningún modelo, lo que coincide con el análisis por grupos de la sección
anterior.

\begin{table}[H]
  \centering
  \begin{tabular}{@{}lrrr@{}}
    \toprule
    \textbf{Grupo} & \textbf{Árbol} & \textbf{Random Forest} & \textbf{XGBoost} \\
    \midrule
    Primer ping & 37\,\% & 45\,\% & 22\,\% \\
    Hoja de ruta & 61\,\% & 48\,\% & 68\,\% \\
    Rastreo previo & 3\,\% & 7\,\% & 10\,\% \\
    \bottomrule
  \end{tabular}
  \caption{Importancia total de cada grupo de atributos.}
  \label{tab:imp-grupo}
\end{table}

El tipo de hoja destaca porque no es un dato cualquiera: cuando está informado los desvíos duran mucho más
(sección siguiente) y falta en la mitad de los eventos, de modo que su ausencia también es información.
Conviene preguntarle a Tecfield qué distingue a esas hojas.

\subsection{Qué variables separan primero el árbol}

La figura~\ref{fig:arbol} muestra las tres primeras preguntas del árbol entrenado con todo el corpus.
En cada nodo figura la cantidad de desvíos que llegan y su duración mediana.

\begin{figure}[H]
  \centering
  \begin{tikzpicture}[
      nodo/.style={draw, rounded corners=2pt, align=center, font=\scriptsize, inner sep=3pt, fill=gray!8},
      lado/.style={font=\scriptsize\itshape},
      level distance=2.1cm,
      level 1/.style={sibling distance=7.6cm},
      level 2/.style={sibling distance=3.8cm},
    ]
    \node[nodo] {\textbf{¿Tiene tipo de hoja?}\\13.632 desvíos, mediana 20,4 min}
      child { node[nodo] {\textbf{Distancia a la ruta $\le$ 2.939 m}\\6.828 desvíos, mediana 40,2 min}
        child { node[nodo] {\textbf{Velocidad $\le$ 33,5}\\4.459, mediana 32,8 min}
          edge from parent node[lado, left] {sí} }
        child { node[nodo] {\textbf{Latitud $\le -38{,}75$}\\2.369, mediana 50,1 min}
          edge from parent node[lado, right] {no} }
        edge from parent node[lado, above left] {sí} }
      child { node[nodo] {\textbf{Distancia a la ruta $\le$ 1.888 m}\\6.804 desvíos, mediana 14,1 min}
        child { node[nodo] {\textbf{Latitud $\le -39{,}16$}\\2.401, mediana 9,6 min}
          edge from parent node[lado, left] {sí} }
        child { node[nodo] {\textbf{Distancia $\le$ 13.959 m}\\4.403, mediana 16,2 min}
          edge from parent node[lado, right] {no} }
        edge from parent node[lado, above right] {no} };
  \end{tikzpicture}
  \caption{Primeros tres niveles del árbol de decisión.}
  \label{fig:arbol}
\end{figure}

Se leen tres cosas:

\begin{itemize}[nosep]
  \item \textbf{La primera pregunta es el tipo de hoja.} Separa dos grupos muy distintos: los 6.828 desvíos
  que lo traen informado tienen una mediana de \textbf{40 minutos} y los 6.804 que no, de \textbf{14
  minutos}. Unos duran casi el triple que otros. Es el mismo atributo que los tres modelos marcan como el
  más importante.
  \item \textbf{La segunda pregunta es la misma en las dos ramas}: qué tan lejos de la ruta estaba el vehículo
  en el primer ping. Con tipo de hoja, los desvíos que se alejaron más de unos 2,9\,km tienen una mediana
  de 50 minutos.
  \item \textbf{El tercer nivel usa la velocidad y la posición}, y el cuarto agrega pasajeros, hoja vigente,
  obra, hora y minutos desde el último GPS. El árbol llega a 8 niveles.
\end{itemize}

Ningún atributo del rastreo previo aparece antes del cuarto nivel, y ahí solo uno (minutos desde el último
GPS), lo que es consistente con su baja importancia.

% ---------------------------------------------------------------- 5
\section{Comparación de todos los modelos}

\begin{table}[H]
  \centering
  \begin{tabular}{@{}lrrrr@{}}
    \toprule
    \textbf{Modelo} & \textbf{MAE (min)} & \textbf{RMSE (min)} & \textbf{Error mediano (min)} & \textbf{Mejora MAE} \\
    \midrule
    Mediana (referencia) & 32,1 & 54,2 & 15,0 & --- \\
    Árbol de decisión & 28,8 & 47,5 & 14,3 & 10\,\% \\
    Random Forest & \textbf{27,0} & \textbf{45,8} & \textbf{13,0} & 16\,\% \\
    XGBoost & 27,3 & 46,1 & 13,1 & 15\,\% \\
    \bottomrule
  \end{tabular}
  \caption{Error de cada modelo con los 23 atributos, en validación cruzada por vehículo. Cuanto más bajo
  el error, mejor. En negrita, el mejor valor de cada métrica.}
  \label{tab:final}
\end{table}

El cuadro~\ref{tab:final} muestra que:

\begin{itemize}[nosep]
  \item \textbf{Todos los modelos le ganan a la referencia} en las tres métricas, y el árbol queda último
  entre los tres modelos en las tres.
  \item \textbf{Random Forest y XGBoost quedan casi empatados.} El Random Forest tiene el menor error en
  las tres métricas, pero la diferencia es chica: 0,3 minutos de MAE, 0,3 de RMSE y 0,1 de error mediano.
  No se midió cuánto varía con otra semilla ni con otras particiones, por lo que no alcanza para decir que
  uno es mejor que otro.
  \item \textbf{El RMSE confirma que el error está en los desvíos largos.} Es alrededor de 1,7 veces el MAE
  en todos los modelos, y la referencia, que no distingue largos de cortos, es la que más lo sufre (54,2 min).
  \item \textbf{El boosting no mejoró al bosque con esta configuración.} Los parámetros de ambos son una
  primera elección sin ajustar, así que esto no descarta que con ajuste XGBoost pueda superarlo.
\end{itemize}

% ---------------------------------------------------------------- 6
\section{Conclusiones}

\subsection{Lo que muestran los resultados}

\textbf{Hay señal.} Con la información de $T_0$ se puede predecir algo de la duración de un desvío: los
tres modelos le ganan a la referencia con vehículos que no vieron al entrenar, y el mejor reduce el MAE de 32,1 a
27,0 minutos y el RMSE de 54,2 a 45,8, cerca de un 16\,\% en ambos.

\textbf{La señal está en el tipo de hoja y en el primer ping.} Los análisis de importancia y la estructura
del árbol apuntan a lo mismo: el tipo de hoja, y luego la distancia a la ruta y la posición. El rastreo
previo casi no aporta.

\textbf{Pero la mejora es modesta.} El mejor modelo se equivoca 27 minutos en promedio y la duración mediana
de un desvío es de 20 minutos. El error sigue siendo grande, sobre todo en los desvíos largos.

\subsection{Limitaciones y pasos siguientes}

\begin{itemize}[nosep]
  \item \textbf{El error real sería mayor.} Los 2.841 desvíos que llegan al tope de cuatro horas no están en
  el corpus y son los más difíciles de predecir, así que el error medido es optimista. Una opción es
  tratarlos como datos censurados con modelos de supervivencia~\cite{hojati}.
  \item \textbf{Los modelos no están ajustados.} Los parámetros de todos son una primera elección;
  falta probar cuáles dan el menor error.
  \item \textbf{Las categorías entran como números}, con un orden que no tienen.
  \item \textbf{Las importancias son del modelo, no causales}, y las de los árboles favorecen a los
  atributos con muchos valores distintos.
  \item \textbf{Hojas vencidas.} El 61\,\% de los puntos de desvío cae fuera de la vigencia de su hoja de
  ruta. Se mantuvieron y se marcaron, pero si no son desvíos válidos habría que sacarlos.
  \item \textbf{Falta rastreo al principio.} El rastreo empieza un mes después que los desvíos, por lo que
  esos eventos no pueden tener rastreo previo.
  \item \textbf{Siguientes pasos:} modelar por separado los desvíos cortos y largos con una etapa previa de
  clasificación~\cite{grigorev}, y reportar el error por segmentos operativos (tipo de vehículo, franja
  horaria, obra) y no solo el global.
\end{itemize}

% ---------------------------------------------------------------- Referencias
\begin{thebibliography}{9}
  \bibitem{konovalenko} A. Konovalenko, L. M. Hvattum y K. A. H. Iversen. Predicting last-mile
  delivery route deviations using machine learning. \emph{Expert Systems With Applications},
  298:129921, 2026.
  \bibitem{grigorev} A. Grigorev, A.-S. Mihaita, S. Lee y F. Chen. Incident duration prediction using
  a bi-level machine learning framework with outlier removal and intra--extra joint optimisation.
  \emph{Transportation Research Part C: Emerging Technologies}, 141:103721, 2022.
  \bibitem{hojati} A. T. Hojati, L. Ferreira, S. Washington y P. Charles. Hazard based models for
  freeway traffic incident duration. \emph{Accident Analysis \& Prevention}, 52:171--181, 2013.
\end{thebibliography}

\end{document}
